\section{Related Work}
\label{sec:related}

Multi-task leaderboards aggregate per-task scores into one ranking. We study two: the HELM MMLU subject tables \citep{liang2023helm,hendrycks2021mmlu} and the Open LLM Leaderboard \citep{ollb2025contents}, whose six task families include BBH \citep{suzgun2023bbh} and are scored with the evaluation harness of \citet{gao2023harness} (Section~\ref{sec:data}).

\paragraph{Fragility of aggregation.}
\citet{dehghani2021lottery} show, on SuperGLUE, VTAB, Long Range Arena and Atari, that which tasks are bundled and how they are combined (mean or median) changes which method comes first. \citet{zhang2024diversity} view a benchmark as an election with tasks as voters, separate cardinal from ordinal aggregation, and measure a trade-off between a benchmark's diversity and its sensitivity to irrelevant changes. \citet{lanctot2023social} also treat tasks as voters and propose social-choice aggregators with robustness properties such as clone-proofness; \citet{singh2025illusion} document how private testing and uneven data access distort the Chatbot Arena ranking. These works vary the \emph{benchmark}, its task or model set or its aggregator. We hold the benchmark fixed and vary the \emph{developer's action}: how many score points, over how many tasks, must be added to one model's scores to make it first under a given rule.

\paragraph{Training on the test task.}
Using evaluation data, or knowledge of the evaluation tasks, at training time raises scores and changes model comparisons \citep{zhou2023cheater,yang2023rephrased,dominguezolmedo2024ttt}. We take such gains as the available lever and train no model: the per-task \emph{cap} on the gain (Section~\ref{sec:model}) stands in for what training can buy, and we vary it, including a cap calibrated on observed fine-tune-minus-base differences.

\paragraph{Strategic post-training.}
The nearest formal treatment of developer incentives is \citet{chen2026incentives}, who model a single benchmark as a Stackelberg game between a designer and competing developers, each choosing an effort that yields a post-effort score, prove that current benchmarks admit no Nash equilibrium among developers, and show that tune-before-test yields an equilibrium ranking by latent quality. Their model has a scalar score per model, so the aggregation of several tasks into one ranking does not arise there, and it studies equilibrium behaviour. Ours is a single developer facing frozen rivals, solved exactly on observed scores, so the two address different quantities.

\paragraph{Bribery and manipulation of scoring rules.}
The closest analogue in computational social choice \citep{brandt2016comsoc} is manipulation and bribery: changing the inputs of an election at minimum cost until a preferred candidate wins \citep{bartholdi1989manipulating,faliszewski2009bribery,faliszewski2010manipulation}, in particular shift and swap bribery, where a bribed voter moves the candidate up its ranking and the cost is the number of positions \citep{elkind2009swap,elkind2020shift}. Those models act on ordinal votes of a population of voters and the question is mostly computational complexity. Ours differs in three respects. The inputs are cardinal scores of a few tasks, the lever is a raise of one model's scores bounded by a per-task cap with cost in score points, and the questions are empirical: how the minimum cost varies across rules, targets and caps on observed leaderboards. The individual solvers are elementary (water-filling for power means, thresholding for the median and the minimum, mixed-integer programs for the ordinal rules and the trimmed mean); what we add is the rule-by-rule comparison of exact minima, the minimum feasible cap $h^\ast$, and the measurements on two leaderboards. We do not claim a new complexity result.

\paragraph{Choice of rules.}
Real leaderboards mostly use the mean: HELM reports a mean win rate over scenarios, the Open LLM Leaderboard the mean of its normalised task scores. We include the geometric and harmonic means and the median as the other common cardinal summaries, the minimum as the limiting case of a worst-task rule, the trimmed mean as a compromise, and Borda, plurality and the Condorcet requirement because the social-choice treatments of benchmarks \citep{zhang2024diversity,lanctot2023social} study ordinal as well as cardinal aggregation of tasks. Plurality and Condorcet are not used by the leaderboards we study; they serve to show how the cost depends on whether scores or ranks are aggregated.
